\section{Workflow thực chiến}

\subsection{Cài đặt môi trường}

Repo hiện có hai cụm phụ thuộc chính.

\subsubsection{Đường XLA}

\begin{lstlisting}[style=rae]
conda create -n rae python=3.10 -y
conda activate rae
pip install uv
uv pip install torch~=2.5.0 torch_xla[tpu]~=2.5.0 torchvision==0.20.1 -f https://storage.googleapis.com/libtpu-releases/index.html
uv pip install timm==0.9.16 accelerate==0.23.0 torchdiffeq==0.2.5 wandb scipy torch-fidelity
uv pip install "numpy<2" transformers einops
\end{lstlisting}

\subsubsection{Phần bổ sung cho JAX / NNX}

\begin{lstlisting}[style=rae]
uv pip install "jax[cuda12]==0.5.1" flax==0.10.4 optax==0.2.4 orbax-checkpoint==0.11.16
uv pip install ml-collections clu absl-py etils huggingface_hub
\end{lstlisting}

Lần chạy JAX đầu tiên sẽ tự bootstrap backend \texttt{diffuse\_nnx} vào thư mục
\texttt{\~{}/.cache/rae\_jax/diffuse\_nnx}.

\subsection{Chuẩn bị model và dữ liệu}

\subsubsection{Model}

\begin{lstlisting}[style=rae]
hf download nyu-visionx/RAE-collections --local-dir models
\end{lstlisting}

\subsubsection{Dữ liệu}

Hầu hết các script dùng dữ liệu dạng \texttt{ImageFolder}. Ví dụ:

\begin{lstlisting}[style=rae]
dataset_root/
  class_a/
    0001.png
  class_b/
    0002.png
\end{lstlisting}

\subsection{Kiểm tra Stage 1 bằng reconstruction}

\subsubsection{Đường XLA / PyTorch}

\begin{lstlisting}[style=rae]
python src/stage1_sample.py \
  --config <config> \
  --image assets/pixabay_cat.png \
  --output recon.png
\end{lstlisting}

\subsubsection{Đường JAX}

\begin{lstlisting}[style=rae]
python3 src_jax/stage1_sample.py \
  --config <config> \
  --image assets/pixabay_cat.png \
  --output recon_jax.png
\end{lstlisting}

Lệnh này phù hợp để xác nhận:

\begin{itemize}[leftmargin=1.5em]
  \item checkpoint decoder Stage 1 có load được hay không,
  \item latent shape có đúng hay không,
  \item preprocessing của encoder có khớp hay không.
\end{itemize}

\subsubsection{Tính latent stat cho Stage 1 trên đường JAX}

Với dataset mới như CelebA, nên tạo một file bootstrap identity stat trước
(\texttt{mean = 0}, \texttt{var = 1}), rồi chạy:

\begin{lstlisting}[style=rae]
python3 src_jax/build_stage1_stats.py \
  --config <stage1_config> \
  --input <train_imagefolder> \
  --output <stage1_stat.pt> \
  --batch-size 16 \
  --set stage_1.params.normalization_stat_path=<bootstrap_identity_stat.pt>
\end{lstlisting}

File \texttt{stage1\_stat.pt} tạo ra vẫn cùng định dạng với repo gốc, nên dùng
lại được cho cả đường \texttt{src/} lẫn \texttt{src\_jax/}.

\subsubsection{Reconstruction cả thư mục trên đường JAX}

\begin{lstlisting}[style=rae]
python3 src_jax/reconstruct_folder.py \
  --config <stage1_config> \
  --input <val_imagefolder> \
  --output-dir recon_jax_dir \
  --batch-size 8
\end{lstlisting}

Lệnh này hữu ích khi muốn export reconstruction set trước khi build FID reference
hoặc khi cần so sánh decoder Stage 1 giữa các checkpoint.

\subsection{Huấn luyện Stage 2 trên TPU bằng TorchXLA}

\begin{lstlisting}[style=rae]
python src/train.py \
  --config configs/stage2/training/ImageNet256/DiTDH-XL_DINOv2-B.yaml \
  --data-path <imagenet_train_split> \
  --results-dir results \
  --image-size 256 \
  --precision bf16
\end{lstlisting}

Trong mỗi optimizer step, train loop thực hiện:

\begin{enumerate}[leftmargin=1.8em]
  \item center crop và đưa ảnh thành tensor;
  \item mã hóa ảnh bằng RAE;
  \item tính transport loss trong latent space;
  \item step optimizer và scheduler trên TPU;
  \item update EMA;
  \item chạy logging, checkpointing, preview sample, validation và FID theo cadence.
\end{enumerate}

\subsection{Huấn luyện Stage 2 bằng JAX / NNX}

\begin{lstlisting}[style=rae]
python3 src_jax/train.py \
  --config configs/stage2/training/ImageNet256/DiTDH-XL_DINOv2-B.yaml \
  --data-path <imagenet_train_root> \
  --results-dir results_jax \
  --precision bf16 \
  --wandb
\end{lstlisting}

Điểm quan trọng:

\begin{itemize}[leftmargin=1.5em]
  \item vẫn dùng cùng file YAML của repo;
  \item hỗ trợ \texttt{--set key=value} để override nhanh từ CLI;
  \item nếu \texttt{stage\_2.ckpt} là file PyTorch \texttt{.pt}, adapter sẽ port
  trọng số sang NNX để khởi tạo;
  \item nếu \texttt{stage\_2.ckpt} là thư mục Orbax, adapter sẽ restore run JAX
  trước đó.
\end{itemize}

\subsection{Wandb logging}

Biến môi trường khuyến nghị:

\begin{lstlisting}[style=rae]
export ENTITY=<wandb_entity>
export PROJECT=<wandb_project>
export WANDB_KEY=<wandb_api_key>
\end{lstlisting}

Đường JAX cũng chấp nhận các tên \texttt{WANDB\_ENTITY} và
\texttt{WANDB\_API\_KEY}. Adapter sẽ bridge các tên legacy ở trên nếu chúng đã
được export sẵn.

\subsection{FID trong train loop}

Đường XLA hỗ trợ online FID trong \texttt{src/train.py}. Ví dụ:

\begin{lstlisting}[style=rae]
eval:
  fid_ref: /path/to/reference_stats.npz
  fid_every: 25000
  fid_num_samples: 4096
  fid_per_proc_batch_size: 4
  fid_batch_size: 128
  fid_device: cpu
  fid_num_threads: 96
\end{lstlisting}

Đường JAX hiện ánh xạ phần FID quan trọng nhất của block \texttt{eval}, nhưng
chưa có parity hoàn toàn cho validation loss như đường XLA.

\subsection{Sampling Stage 2}

\subsubsection{JAX một lệnh}

\begin{lstlisting}[style=rae]
python3 src_jax/sample.py \
  --config configs/stage2/sampling/ImageNet256/DiTDHXL-DINOv2-B_AG.yaml \
  --class-labels 207,360 \
  --output sample_jax.png
\end{lstlisting}

\subsubsection{JAX distributed}

\begin{lstlisting}[style=rae]
python3 src_jax/sample_ddp.py \
  --config configs/stage2/sampling/ImageNet256/DiTDHXL-DINOv2-B.yaml \
  --sample-dir samples_jax \
  --num-samples 50000 \
  --label-sampling equal \
  --fid-ref /path/to/reference_stats.npz
\end{lstlisting}

\subsection{Đẩy artifact lên Hugging Face}

\subsubsection{Lệnh độc lập}

\begin{lstlisting}[style=rae]
python3 src_jax/push_hf.py \
  --path results_jax/<run_name> \
  --repo-id <hf_user_or_org>/<repo_name>
\end{lstlisting}

\subsubsection{Tích hợp ngay trong train}

\begin{lstlisting}[style=rae]
python3 src_jax/train.py \
  --config <config> \
  --data-path <train_root> \
  --hf-repo-id <hf_user_or_org>/<repo_name>
\end{lstlisting}

\subsection{Notebook Kaggle cho CelebA}

Repo hiện có notebook \texttt{raes-jax-celeba-kaggle.ipynb} để bám đúng flow
Kaggle:

\begin{itemize}[leftmargin=1.5em]
  \item clone repo và checkout branch \texttt{jax};
  \item tạo \texttt{.venv} riêng bằng \texttt{uv};
  \item chạy các bước phụ thuộc package qua \texttt{uv run};
  \item chuyển CelebA thành \texttt{ImageFolder} \texttt{256x256};
  \item tạo bootstrap identity stat, build latent stat cho Stage 1;
  \item export reconstruction validation set, build \texttt{fid\_ref};
  \item ghi sẵn config CelebA cho \texttt{src\_jax/train.py}.
\end{itemize}

Nếu chạy trên Kaggle \texttt{TPU v5e-8}, dùng notebook
\texttt{raes-jax-celeba-kaggle-tpuv5e8.ipynb}. Bản này đổi sang đường cài
\texttt{jax[tpu]}, giữ các bước phụ thuộc package trong \texttt{.venv} riêng
qua \texttt{uv run}, tự xóa cờ executable-stack mà Kaggle có thể chặn trên
\texttt{jaxlib}, chạy bước xác nhận TPU trong một tiến trình Python mới, giữ
Stage 1 và Stage 2 ở TPU, còn FID và thống kê tham chiếu sẽ chạy trên host CPU
nhiều luồng.
